% ===========================================================================
%  Page de garde
% ===========================================================================
\begin{titlepage}
\newgeometry{margin=0cm}
\thispagestyle{empty}
\begin{tikzpicture}[remember picture, overlay]
  % ---------------------------------------------------------- fond marine
  \fill[marine] (current page.north west) rectangle ([yshift=-11.2cm]current page.north east);
  \fill[ocean]  ([yshift=-11.2cm]current page.north west)
                rectangle ([yshift=-11.5cm]current page.north east);

  % ------------------------------------------------- plateau stylisé (décor)
  \begin{scope}[shift={([xshift=-8.3cm, yshift=-9.6cm]current page.north east)},
                scale=0.62, opacity=0.9]
    \draw[step=1, line width=0.3pt, draw=white!35!marine] (0,0) grid (10,10);
    \draw[line width=0.9pt, draw=ocean] (0,0) rectangle (10,10);
    % navires : 1x4, 2x3, 3x2, 4x1, sans contact
    \foreach \x/\y/\w/\h in {1/8/4/1, 7/5/1/3, 2/2/3/1, 0/5/2/1, 4/5/1/2, 8/1/2/1,
                             6/9/1/1, 9/9/1/1, 0/0/1/1, 6/3/1/1}
      \fill[ocean, rounded corners=1.5pt]
           (\x+0.12,\y+0.12) rectangle (\x+\w-0.12,\y+\h-0.12);
    % coups au but
    \foreach \x/\y in {1/8, 2/8, 7/6, 4/5}
      \fill[corail] (\x+0.5,\y+0.5) circle (0.2);
    % coups dans l'eau
    \foreach \x/\y in {3/6, 5/8, 8/3, 1/4, 6/1, 9/6, 3/0, 5/3}
      \fill[white!55!marine] (\x+0.5,\y+0.5) circle (0.1);
  \end{scope}

  % ------------------------------------------------------------- en-tête
  \node[anchor=north west, text=white, font=\sffamily\small, align=left, inner sep=0pt]
    at ([xshift=2.6cm, yshift=-2.2cm]current page.north west)
    {\textbf{Sup Galilée} --- Institut Galilée\\[2pt]
     Université Sorbonne Paris Nord};

  % --------------------------------------------------------------- titre
  \node[anchor=north west, text=white, align=left, inner sep=0pt, text width=9.2cm]
    at ([xshift=2.6cm, yshift=-4.6cm]current page.north west)
    {{\fontsize{44}{48}\selectfont\sffamily\bfseries Bataille\\[4pt] Navale}\\[18pt]
     {\Large\sffamily Du travail pratique\\[2pt] à l'architecture logicielle}};

  \node[anchor=north west, text=ocean!30, font=\sffamily, inner sep=0pt]
    at ([xshift=2.6cm, yshift=-10.2cm]current page.north west)
    {Rapport de projet \textbullet\ version 2.1};

  % ---------------------------------------------------------------- corps
  \node[anchor=north west, align=left, inner sep=0pt, text width=15.5cm]
    at ([xshift=2.6cm, yshift=-13.4cm]current page.north west)
    {{\sffamily\color{ardoise}\small Unité d'enseignement}\\[4pt]
     {\Large\sffamily\bfseries\color{marine} Initiation aux interfaces graphiques}\\[3pt]
     {\sffamily Cycle préparatoire ingénieur intégré, 2\textsuperscript{e} année
      \textbullet\ option, semestre~4}\\[3pt]
     {\sffamily Enseignante : Sophie \textsc{Toulouse}}\\[1.3cm]
     %
     {\sffamily\color{ardoise}\small Auteurs}\\[4pt]
     {\Large\sffamily\bfseries\color{marine} Dmytro \textsc{Palahin}
      \quad\textcolor{ocean}{\textbullet}\quad Ivanna \textsc{Makam}}\\[6pt]
     {\sffamily\small\color{ardoise} Projet réalisé en binôme en 2023 (version~1).\\
      Refonte de l'architecture et rédaction du présent rapport (version~2) :
      Dmytro \textsc{Palahin}, 2026.}\\[1.3cm]
     %
     {\sffamily\color{ardoise}\small Technologies}\\[4pt]
     {\sffamily C11 \textbullet\ GTK~3 \textbullet\ CMake \textbullet\ GitHub Actions}};

  % ----------------------------------------------------------------- pied
  \draw[ocean, line width=1.2pt]
    ([xshift=2.6cm, yshift=2.6cm]current page.south west) -- ++(2.2cm,0);
  \node[anchor=south west, font=\sffamily\small, text=ardoise, align=left, inner sep=0pt]
    at ([xshift=2.6cm, yshift=1.6cm]current page.south west)
    {Code source : \url{https://github.com/DmytroPalahin/Battleship}\\[2pt]
     Novembre 2026};
\end{tikzpicture}
\restoregeometry
\end{titlepage}
